% page 345 du fac-simile (image p-344 du scan)
% bloc 160.0 x 237.7 mm ; marge gauche 42.0 mm
\pgc{20.320mm}{0.000mm}{
\l{\cel{52}337.}
\l{}
\l{\sou{libraco}. (\sou{astron.}) Ocilo videbla di astro cirkum sua axo.}
\l{\cel{6}DEFI.}
\l{}
\l{\sou{lico}. Tereno +kluza per barili, en qua luktas la personi qui}
\l{\cel{3}partoprenas turniri, konkursi, e c. - EFIS.}
\l{}
\l{liceo. (en\cel{1}\sou{Francia}). Edifico en qua agesas la doco sekundara,}
\l{\cel{3}direktata da la stato (kontraste kun \textquotedbl{}kolegio\textquotedbl{}). - FI.}
\l{}
\l{\sou{licensar.(trans.})\cel{2}(\sou{aludante stat-administrantaro)}. Koncesar}
\l{\cel{3}partikulare la yuro pri debitar ula vari konsumebla, o pri}
\l{\cel{3}peskar, chasar, e c., en ta o ca regiono. DEFIRS.}
\l{}
\l{\sou{licenco}.I(En\cel{1}\sou{Francia}). Diplomo universitatala, inter bachele-}
\l{\cel{3}reso-diplomo e doktoreso-diplomo, e qua yurizas pri docar}
\l{\cel{3}en licei o kolegii, o pri pledar, e c. - II.\cel{1}\sou{Licenco}}
\l{\cel{3}\sou{poeziala}\cel{1}: su-liberigo, da poeto, pri la strikta reguli}
\l{\cel{3}prozodiala. - DEFIRS.}
\l{}
\l{\sou{lido}. Parto movebla quan onu pozas sur (o qua retrovenas}
\l{\cel{3}ad) la aperturo di vazo, di kofro, e c., por kovrar o klo-}
\l{\cel{3}zar lu. - E.}
\l{}
\l{\sou{lienteri}o.(\sou{patol.}) Diareo en qua la nutrivi esas forjetita}
\l{\cel{3}mi-digestita. -}
\l{}
\l{\sou{lietnanto}. Oficiro di qua la\cel{1}\sou{gra}do esas nemediate sub olta}
\l{\cel{3}di kapitano (en armeo).DEFIRS.}
\l{}
\l{\sou{lifto}. Mashino qua elevas la personi a la etaji diversa di}
\l{\cel{3}edifico. - EF.}
\l{}
\l{ligar. (\sou{trans}.)I. Cirkondar la parti di objekto o di plura}
\l{\cel{3}objekti, por interjuntar li, per ulo flexebla, ulo longa}
\l{\cel{3}(kordo, rubando, e c.). - II.JUntar (plura kozi) per relato}
\l{\cel{3}di sucedo, konsequo, e c. - III. Unionar (plura personi)}
\l{\cel{3}per relati di socio, di amikeso, e c. - EFIRS.}
\l{}
\l{ligamento. (\sou{anat.)}\cel{2}Fasko de tisuo fibroza densa, poke exten-}
\l{\cel{3}sebla, qua ligas osti, kartilagi, muskuli. - DEFIS.}
\l{}
\l{ligno. Parto dura, ek la tisuo fibroza di arboro, di qua la}
\l{\cel{3}pezo specifika e la dikeso-grado varias segun la speci e la}
\l{\cel{3}evo. - EFIS.}
\l{}
\l{lignino.Substanco kemiala qua impregnas la elementi di ligno,}
\l{\cel{3}e qua donas ad ica lua grado di fermeso, di solideso. DEFIS.}
\l{}
\l{lignito. Karbono fosila qua konservas en su, plu o min multe,}
\l{\cel{3}traci di konstituciono vejetala. - DEFIRS.}
\l{}
\l{liguo. Asociitaro formacita interne di stato por igar triumfar}
\l{\cel{3}ula interesti politikala, religiala. - EF.}
}
